% === Introduction ===
\section{Introduction}
\label{sec:introduction}

Modern cloud services emit a relentless stream of alerts, and on-call engineers
must decide, often within minutes, which signal a genuine incident and which are
noise~\cite{he2021survey,ghosh2022incidents}. The decision is consequential both
ways: a missed incident becomes an outage, while a flood of low-value pages
erodes trust in the alerting system. What makes an on-call engineer effective is
rarely raw signal-reading; it is \emph{memory}. Most incidents rhyme with one the
team has resolved before, and the fastest path from symptom to fix is to
recognize ``the message-broker outage from last quarter'' and reuse the
diagnosis that worked.

This has motivated a large literature on automated incident triage and
AIOps~\cite{he2021survey,chen2019triage,ghosh2022incidents}, increasingly powered
by large language models~\cite{ahmed2023rootcause,chen2024rcacopilot}. That
literature is evaluated in a strikingly uniform way: faults are scripted into a
demonstration microservice application, telemetry is collected, and a system is
scored on whether it can tell a faulty window from a healthy
one~\cite{chen2025aiopslab}. Reported numbers are high and still improving.

This paper asks whether that is the right thing to measure.

\smallskip
\noindent\textbf{The problem.} We find that the dominant benchmark target is
saturated by construction and does not survive contact with real incidents.
Deciding whether a window warrants a ticket is almost perfectly solvable on
injected faults --- an active fault differs from a healthy baseline by
design --- yet only weakly predictable from the text of real, engineer-written
tickets. On our corpora the same task moves from ROC area $1.00$ to
$0.57$--$0.70$. A strong score on fault-injection data therefore certifies the
dataset, not the method.

Worse, we find that several evaluation choices which are routinely left
implicit change not just the magnitude of a result but its \emph{direction}.
Whether retrieval fusion helps or hurts depends on the size of the corpus it is
measured on. Whether a reported margin is significant depends on whether the
training seed is treated as a source of variance. Whether a method looks best
depends on which of two defensible relevance relations is used. We know these
because each of them reversed or materially altered one of our own conclusions.

\smallskip
\noindent\textbf{What we do.} We run a measurement study across three corpora of
increasing realism: two newly released microservice-telemetry datasets that pair
multi-signal telemetry with incident tickets and per-window labels (Online
Boutique and OTel Demo)\footnote{\url{https://zenodo.org/records/21096887}} and
$38{,}642$ real Apache Jira incidents across 24 projects (World of
Logs~\cite{wang2026worldoflogs}). Throughout we use \sys{}, a retrieval system
that fuses dense, lexical, and knowledge-graph matching behind a cost-aware
controller, as the \emph{instrument}: a strong, realistic baseline whose
behaviour exposes the measurement problems we are studying. Several of our findings
are unflattering to it, which is the point --- the object of study is the
evaluation, not the system.

\smallskip
\noindent\textbf{Contributions.}
\begin{itemize}
  \item \textbf{Evidence that the dominant target does not transfer.}
  Ticket-worthiness classification is saturated on injected faults and only
  weakly learnable on real incidents, while incident-memory retrieval is
  measurable identically on both and does transfer. We also show that our own
  earlier reading of this result was an artifact of a degenerate feature matrix,
  and how a one-line check catches it.

  \item \textbf{A quantification of how much protocol choices move results.}
  The training seed accounts for up to $28\%$ of variance in a headline metric;
  GPU fine-tuning does not reproduce at a fixed seed; the choice of relevance
  relation moves \hit{5} by $0.23$ and reorders methods; and retrieval-fusion
  conclusions \emph{reverse} between a $347$-item and a $38{,}642$-item memory.

  \item \textbf{Two released telemetry-to-ticket datasets} pairing logs,
  metrics, traces, and cluster events with incident tickets and per-window
  labels (26 and 52 fault scenarios), to our knowledge the first public corpora
  coupling telemetry and tickets at this granularity, alongside all code,
  baselines, ablations, and significance
  tests.\footnote{\url{https://zenodo.org/records/21096898}}

  \item \textbf{A protocol for evaluating incident triage}
  (\S\ref{sec:recommendations}), each item tied to a specific way our own
  results would have misled without it.
\end{itemize}

\noindent We report findings that are unfavourable to the instrument as
readily as favourable ones: no evidence-gathering tool in \sys{} has a
measurable benefit, its escalation gates fire roughly $16\times$ more often than
necessary, and its retrieval fusion improves top-five recall while
significantly degrading top-one precision. These are the kinds of result a
system paper is tempted to omit, and the kinds a measurement paper exists to
report.
